%!TEX root = TQFTnotes.tex
%%%%%%%%%%%%%%%%%%%%%%%%%%%%%%
\chapter{Separable Frobenius algebras}\label{c:separable-frobenius}

In this chapter, we study duality in the 2-category $\Alg$ of algebras and
bimodules, which was introduced in \exref{x:alg}. It turns out that the maximal
subcategory with duals in $\Alg$ is the category of separable algebras (which in
characteristic zero is the same as finite-dimensional semisimple algebras); if
we also want to trivialize the Serre functor (in other words, if we want to have
adjunction between evaluation and coevaluation 1-morphisms), we arrive at the
notion of a separable symmetric Frobenius algebra. These results will be heavily
used later, when we classify extended 2-dimensional TQFTs.

To simplify the exposition, throughout this chapter we assume that $\kk$ is a
field of characteristic zero.

%%%%%%%%%%%%%%%%%%%%%%%%%%%%%%%
\section{Duality in $\Alg$}\label{s:alg-duals}
Our goal is analyzing duality in the 2-category $\Alg$, looking at duality
both for objects and for 1-morphisms. Recall that objects of this category are
associative algebras, and 1-morphisms $A\to B$ are $B$-$A$-bimodules.


Let us begin by looking at duality for objects. It turns out that each $A\in
\Alg$ has a dual. Namely,\index{opposite!algebra} for an algebra $A$, denote by
$A^\op$ the \emph{opposite algebra}, which coincides with $A$ as a vector space
but has the opposite multiplication. Then any $A$-bimodule can also be considered
as a left module over the algebra $A\otimes A^\op$, or as a right module over
$A\otimes A^\op$. In particular, $A$ itself is a left module over $A\otimes
A^\op$, and also a right module over $A\otimes
A^\op$.


\begin{theorem}\label{t:alg-duals1}
  Let $A\in \Alg$. Define the evaluation and coevaluation morphisms
  $A^\op \otimes A\to \kk$, $\kk \to A\otimes A^\op$
  to be $A$, considered either as a left or a right $A\otimes A^\op$ module:
  \begin{equation}\label{e:alg-duals1}
    \begin{aligned}
      \ev &= A_{A\otimes A^\op} \colon A\otimes A^\op\to \kk\\
      \coev &= {}_{A\otimes A^\op} A \colon \kk \to A\otimes A^\op.
    \end{aligned}
  \end{equation}
  Then $(A, A^\op, \ev, \coev, \al,\be)$ is a coherent dual pair in $\Alg$,
  with isomorphisms $\al, \be\colon A\otimes_A A \isoto A$ given by
  multiplication.
\end{theorem}
In other words, any algebra $A$ has a dual, namely $A^\op$.
Note that it has nothing to do with the dual vector space, and
it works for finite-dimensional as well as infinite-dimensional algebras.

Let us now look at duality for 1-morphisms, i.e., bimodules.
Recall that a module $P$ over an algebra $A$ is called
\emph{projective}\index{projective module}
if it is a direct summand of a free module.

\begin{theorem}\label{t:alg-adjoint}
  Let ${}_BM_A$ be a $B$-$A$-bimodule, considered as a 1-morphism
  $A\to B$ in $\Alg$.
  Then the following are equivalent:
  \begin{enumerate}
    \item $M$ has a right adjoint.
    \item There exists an $A$-$B$-bimodule $N$, an element
      $e=\sum_{i=1}^k n_i\otimes m_i\in N\otimes M$,
      and a morphism of $B$-bimodules $(\, ,\,)\colon M\otimes_A N\to B$
      such that $e$ is central: for any $x\in A$,
      $$
        \sum (xn_i)\otimes m_i=\sum n_i\otimes (m_i x)
      $$
      and for any $m\in M$, $n\in N$ we have
      \begin{equation}\label{e:alg-adjoint-snake}
        \begin{aligned}
          m&=\sum (m, n_i) m_i,\\
          n&=\sum n_i (m_i, n).
        \end{aligned}
      \end{equation}
    \item $M$ is a finitely generated projective $B$-module.
  \end{enumerate}
  In this case, its right adjoint is isomorphic to the module
  $N=M^\vee = \Hom_B (M, B)$ \textup{(}note that it has a
  canonical $A$-$B$-bimodule structure\textup{)}.
\end{theorem}
\begin{proof}
  Equivalence (1)$\iff$ (2) is obvious from the definition; the form
  $(\, , \,)$ plays the role of evaluation morphism, and $e=\eta (1)$,
  where $\eta\colon A\to N\otimes M$ is the coevaluation.

  To prove (2)$\implies$ (3), assume that we are given such an
  $e=\sum_{i=1}^k n_i\otimes m_i$. Consider the free
  $B$-module $F$ of rank $k$: $F=\bigoplus B x_i$. Define morphisms
  of $B$-modules $s\colon M\to F$, $p\colon F\to M$ by
  \begin{align*}
  s(m)&=\sum (m, n_i)x_i\\
  p(x_i)&=m_i.
  \end{align*}
  Then it follows from (2) that $p\circ s =\id_{M}$, which implies that
  $M$ is a direct summand in $F$, so $M$ is projective of finite rank.
  Thus, (2)$\implies$ (3).

  Proof of (3)$\implies$ (2) is similar and is left to the reader.

  Finally, to show that in this case, the right adjoint $N$ is isomorphic
  to $M^\vee$, we construct morphisms $N\to M^\vee$, $M^\vee\to N$ as follows:
  \begin{align*}
    N&\to M^\vee\colon n\mapsto (-, n)\\
    M^\vee&\to N\colon f\mapsto \sum n_i f(m_i)
  \end{align*}
  We leave it to the reader to verify that the so-defined maps are indeed morphisms
  of right $A$-modules, and that they are mutually inverse.
\end{proof}

Similarly, $M$ has a left adjoint iff it is finitely generated and projective
as a right $A$-module.

%%%%%%%%%%%%%%%
\section{Separable algebras}\label{s:separable}

\thref{t:alg-adjoint} shows that in order for the coevaluation morphism $\kk
\to A^\op\otimes A$ (which is given by $A$ considered as a left $A^\op \otimes
A$ module) to have a right adjoint in $\Alg$, it is necessary that $A$ is a
projective $A^\op\otimes A$ module. This naturally leads us to the notion of a
\emph{separable} algebra. We give a brief review of this notion here, referring
the reader to \ocite{pierce}*{Chapter 10} for details.


\begin{lemma}\label{l:separable}\index{algebra!separable}
  Let $A$ be an associative algebra over a field $\kk$.
  Then the following are equivalent:
  \begin{enumerate}
    \item $A$ is projective as a left module over $A^e=A\otimes A^\op$.
    \item The multiplication map $m\colon A\otimes A\to A$,
      considered as a morphism of $A$-bimodules
      ${}_AA\otimes A_A \to {}_A A_A$,
      has a one-sided inverse: there exists a morphism of $A$-bimodules
      $\De\colon A\to A\otimes A$ such that $m\circ \De=\id_A$.
   \item There exists an element $E=\sum x_i\otimes y_i\in A\otimes A$ which
     is central: $\sum ax_i \otimes y_i=\sum x_i \otimes y_ia$ and satisfies
     $m(E)=\sum x_iy_i=1$.
   \end{enumerate}
\end{lemma}
\begin{proof}
  (1)$\implies$(2) is a standard fact of basic algebra: if $P$ is a projective
  module, then any surjective morphism $X\to P\to 0$ has a one-sided inverse.

  Similarly, (2)$\implies$(1) is immediate from the definition: if such a
  morphism $\De$ exists, then $A$ is a direct summand in
  the free $A^e$-module $A\otimes A\simeq A^e$.

  Finally, (3) is just a rewriting of (2).
\end{proof}
\begin{definition}\label{d:separable}
  An algebra satisfying one of the equivalent conditions of \leref{l:separable}
  is called \emph{separable}.
\end{definition}
An element $E\in A\otimes A$ as in part (3) of the lemma above is called a
\emph{separability idempotent}\index{separability idempotent}. This terminology
is justified by the following exercise.
\begin{exercise}\label{ec:separability-idempotent}
  Let $A$ be a separable algebra with separability idempotent $E$.
  Then $E$, considered as an element in the algebra $A\otimes A^\op$,
  satisfies $E^2=E$.
\end{exercise}


Note that the separability idempotent is not unique, nor does it have
to be symmetric, as is illustrated by the following example.
\begin{example}\label{x:matrix}
  Let $A=\Mat_n(\kk)$ be the matrix algebra. Choose an index $k$ and define
  $e_k=\sum_i E_{ik}\otimes E_{ki}$, where $E_{ij}$ are matrix units.
  Then each of $e_k$ is a separability idempotent, so $A$ is separable.

  We can also take as the separability idempotent the element
  $$
    E=\tfrac{1}{n}\sum E_{ij}\otimes E_{ji}
  $$
  which has the added benefit of being symmetric.
\end{example}
\begin{remark}
  It can be shown that a \emph{symmetric} separability idempotent, if it exists, is
  unique, see \ocite{aguiar}*{Theorem 3.1}.
\end{remark}


One of the main results of this theory is the following.
\begin{theorem}\label{t:separable-ss}
  If $\kk$ has characteristic zero, an algebra $A$ over $\kk$ is
  separable if and only if it is finite-dimensional and semisimple.
  In this case, for any field extension $\kk\subset F$, the algebra
  $A_F=A\otimes_\kk F$ is also semisimple over $F$.
\end{theorem}
We skip the proof of this result, referring the reader to
\ocite{pierce}*{Section 10.7}.

\begin{remark}
  In fact, \thref{t:separable-ss} can be generalized: it works for any perfect
  field, i.e., a field such that any finite extension of it is separable in the
  sense used in Galois theory. This includes not only fields of characteristic
  zero, but also finite fields.
\end{remark}




Since any module over a semisimple algebra is projective,
the results of the previous section immediately imply the following theorem.
\begin{theorem}\label{t:ssalg}
  Let $\ssAlg\subset \Alg$ be the category whose objects are separable
  algebras, and 1-morphisms are finite-dimensional
  bimodules. Then $\ssAlg$ is a category with duals as defined in
  \deref{d:category-duals}.
\end{theorem}

\begin{exercise}\label{ec:serre-separable}
  Show that for a separable algebra $A$, the Serre automorphism as defined in
  \deref{d:serre} is given by $A^*$ (dual vector space), considered as
  an $A$-bimodule.
\end{exercise}




%%%%%%%%%%%%%%%%%%%%%%%%%%%%%%%
\section{Separable symmetric Frobenius algebras}\label{s:ssfrob}
The results of the previous section show that fully dualizable objects in $\Alg$ are
separable algebras. However, for future use, let us also consider this question:
what are fully dualizable algebras for which the evaluation $\ev_A$ and
coevaluation $\coev_A$ are adjoints of each other?


One way to answer that is to use the notion of Serre functor, introduced in
\seref{s:serre}. By the results of that section, if $\ev$ has left and right
duals $\ev^L$, $\ev^R$, then they coincide with the coevaluation $\coev$ twisted
by the Serre automorphism or its inverse. Thus, to construct adjunctions
$\ev\dashv\coev$, $\coev\dashv\ev$ we need to construct an isomorphism
$S_A\simeq 1_A$. Since by \ecref{ec:serre-separable} the Serre automorphism for
an algebra $A$ is given by $A^*$ considered as an $A$-bimodule, we see that
trivialization of the Serre automorphism is the same as a choice of an
isomorphism of $A$-bimodules $A\simeq A^*$.

It turns out that the answer is related to Frobenius algebras. Recall that a
symmetric Frobenius algebra is a finite-dimensional algebra $A$ together with a
linear map $\eps\colon A\to\kk$ such that the pairing
$$
  A\otimes A\to \kk\colon
  a\otimes b \mapsto (a,b)=\eps(ab)
$$
is a symmetric non-degenerate bilinear form, see \deref{d:frobenius2},
\deref{d:frobenius_symmetric}.




\begin{exercise}\label{ec:fd-algebra-separable}
  Show that for a finite-dimensional algebra $A$, constructing an isomorphism of
  $A$-bimodules $A\simeq A^*$ is equivalent to defining on $A$ a structure of
  a symmetric Frobenius algebra (compare with \ecref{ec:dual-module-frob}).
\end{exercise}


However, it is useful to also give a more direct construction, relating
adjunctions $\ev\dashv\coev$, $\coev\dashv\ev$ defined by \eqref{e:alg-duals1}
with the Frobenius algebra structure.

\begin{theorem}\label{t:adjunction-frob}
  Let $A$ be an algebra considered as an object in $\Alg$. Then defining
  adjunction data for $\ev_A\dashv\coev_A$ is equivalent to defining on $A$ a
  structure of a symmetric Frobenius algebra.
\end{theorem}
In particular, this theorem shows that such an adjunction is only possible if
$A$ is finite-dimensional.
\begin{proof}

  By definition, adjunction data consists of the unit and the counit of adjunction, i.e.,
  2-morphisms $\eps \colon \ev_A\circ \coev_A \to 1_\kk$, $\eta \colon
  1_{A\otimes A^\op}\to \coev_A\circ \ev_A$, where $\ev_A\colon A\otimes
  A^\op\to\kk$, $\coev_A \colon \kk \to A\otimes A^\op$ are evaluation and
  coevaluation 1-morphisms.

  Recalling that $\ev_A$ is given by $A$ considered as a right $A^e$ module, and
  $\coev_A$ is given by $A$ considered as a left $A^e$ module, we see that
  composition of 1-morphisms $\ev_A \circ \coev_A$,
  considered as a $\kk$-bimodule, in other words a vector space, is given by
  $$
  \ev_A \circ \coev_A \simeq A\otimes_{A^e} A\simeq A/[A,A]
  $$
  where $[A,A]$ is the subspace in $A$ generated by $xy-yx, x,y\in A$.
  Note that $A/[A,A]$ is just a vector space; it has no natural algebra
  structure (unlike the situation with Lie algebras, where $L/[L,L]$ is
  a Lie algebra), since $[A,A]$ is not an ideal in $A$.


  Thus, we see that the counit of adjunction
  $\eps\colon \ev_A \circ \coev_A\to 1_\kk$
  is a linear map
  $$
    \eps\colon A/[A,A]\to \kk
  $$
  or equivalently, a map $\eps \colon A\to \kk$ such that $\eps(xy)=\eps(yx)$.

  Similarly, the unit of adjunction should be a morphism of bimodules
  ${}_{A^e} A^e_{A^e} \to {}_{A^e}A\otimes A_{A^e}$. Unpacking this definition,
  we see that $\eta$ is a map of modules over 4 copies of $A$:
  \begin{equation}\label{e:eta-ssfrob}
  \eta\colon {}_{A_1} A_{A_2}\otimes {}_{A^\op_3}A_{A^\op_4}
      \to   {}_{A_1} A_{A_3}\otimes {}_{A_4}A_{A_2}
  \end{equation}
  where indices are used to indicate copies of $A$. It is immediate that such a
  morphism is uniquely determined by $e=\eta(1\otimes 1)$, which must be
  \emph{bicentral}\index{bicentral element}: if $e=\sum x_i\otimes x^i$, then
  \begin{equation}\label{e:bicentral}
     \begin{aligned}
       \sum xx_i\otimes x^i&=\sum x_i\otimes x^ix,\\
       \sum x_iy\otimes x^i&=\sum x_i\otimes yx^i.
     \end{aligned}
   \end{equation}
   We leave it to the reader to check that the adjunction conditions on $\eta$,
   $\eps$ are equivalent to the conditions that for any $x\in A$, we have
   $$
     x= \sum (x, x_i)x^i = \sum (x, x^i) x_i,
   $$
   where $(a,b)=\eps(ab)$ (compare with the proof of \thref{t:frobenius1}).
   This means that $(\, ,\,)$ is a non-degenerate pairing and $x_i, x^i$
   are dual bases with respect to this pairing, so $A$ is a symmetric Frobenius
   algebra.
   Conversely, if $A$ is a symmetric Frobenius algebra, we can define
   $e=\sum x_i \otimes x^i$ where $x_i$, $x^i$ are dual bases with respect to
   the Frobenius pairing. By \thref{t:frobenius2}, the so-defined $e$ is central;
   since it is symmetric, it is also bicentral and thus defines a morphism
   $\eta$ as in \eqref{e:eta-ssfrob}.

\end{proof}

\begin{remark}
  If you have difficulty keeping track of the many copies of $A$ and in which
  order they appear, you are not alone. Fortunately, it will become clearer
  once we relate it to geometry: see \seref{s:extended-alg}.
\end{remark}


So far we have discussed adjunction $\ev\dashv \coev$. What about the other
adjunction, $\coev\dashv \ev$?

First, by \thref{t:alg-adjoint}, in order for $\coev={}_{A^e}A$ to have a right
adjoint it is necessary that $A$ is projective over $A^e$, that is, that $A$ is
separable. It turns out that this is also sufficient --- if we assume we
already have adjunction $\ev\dashv\coev$.


\begin{theorem}\label{t:adjunction-frob2}
  Let $A$ be an algebra together with adjunction data $\ev_A\dashv\coev_A$,
  $\coev_A\dashv\ev_A$. Then $A$ is separable and has a canonical structure of a
  symmetric Frobenius algebra, with the counit $\eps\colon A/[A,A]\to \kk$ given
  by the counit of adjunction $\ev\dashv\coev$.

  Conversely, let $A$ be a separable symmetric Frobenius algebra. Then we can
  define adjunction data $\ev_A\dashv\coev_A$, $\coev_A\dashv\ev_A$.
\end{theorem}

To prove this theorem, we will need two preliminary results.

\begin{theorem}\label{t:adjunction-frob3}
  For an algebra $A$, adjunction data for adjunction $\coev_A\dashv\ev_A$ is
  equivalent to an element $C\in A/[A,A]$ and a bicentral element
  $$
    \tilde e=\sum\tilde x_i\otimes \tilde x^i\in A\otimes A
  $$
  such that $\sum \tilde x_i C \tilde x^i=1$.
\end{theorem}
Note that bicentrality of $\tilde e$ implies that $\sum \tilde x_i C \tilde
x^i$ only depends on the class of $C$ in $A/[A,A]$.

The proof of this theorem is similar to the proof of \thref{t:adjunction-frob};
the element $C$ corresponds to the unit $\kk \to \ev_A \circ \coev_A$, and $\tilde
e$ corresponds to the counit $\coev_A\circ \ev_A \to 1_{A^e}$. We leave the
details to the reader.



Before proving the next theorem, we need some notation. Let $A$ be a symmetric
Frobenius algebra and let $x_i, x^i$ be dual bases with respect to the Frobenius
pairing. Consider the map $f\colon A\to A$ given by $x\mapsto \sum x_i xx^i$.
Then bicentrality of $e$ implies that for any $x$, $f(x)\in z(A)$, where $z(A)$
is the center of $A$. Similarly, it also shows that $f(xy)=f(yx)$, so we can
consider $f$ as a map
\begin{equation}\label{e:coinv-inv}
  \begin{aligned}
    A/[A,A]&\to z(A)\\
    [x]&\mapsto \sum x_ixx^i.
  \end{aligned}
\end{equation}



\begin{theorem}\label{t:separable-frobenius}
  Let $A$ be a symmetric Frobenius algebra over a field $\kk$ of characteristic
  zero. Let $x_i, x^i$ be dual bases in $A$ with respect to the Frobenius
  pairing. Then the following are equivalent:
  \begin{enumerate}
    \item $A$ is separable.
    \item $A$ is semisimple.
    \item There exists $[C]\in A/[A,A]$ such that $\sum x_iCx^i=1$.
    \item The Euler element \textup{(}see \eqref{e:euler}\textup{)}
      \begin{equation}\label{e:euler2}\index{Euler element}
        w=\sum x_ix^i
      \end{equation}
      is invertible.
  \end{enumerate}
\end{theorem}
\begin{proof}
  Equivalence of (1) and (2) is \thref{t:separable-ss}
  (note that a Frobenius algebra is automatically finite-dimensional).

  (3)$\implies$(1) is immediate: we can take $\sum x_i C\otimes x^i$ as
  separability idempotent.

  (4)$\implies$(3) is also clear: if $w=\sum x_i x^i$ is invertible, then its
  inverse $w^{-1}$ must be central, and $\sum x_i w^{-1} x^i=1$.

  It remains to prove (1)$\implies$(4).

  Assume that $A$ is separable; then $A$ is semisimple and moreover, if
  $\bar \kk$ is the algebraic closure of $\kk$, then
  $A_{\bar \kk}=A\otimes_\kk \bar\kk$ is semisimple. By
  \leref{l:frobenius-ss}, $A_{\bar\kk}\simeq\bigoplus \Mat_{n_i}(\bar \kk)$,
  and $\eps(a)=\sum \la_i \tr(a_i)$.
  An easy explicit computation then shows that the element $e$ is given by
  $$
    e = \sum \la_i^{-1}E^i_{kl}\otimes E^i_{lk}
  $$
  where $E^i_{kl}$ are matrix units in $\Mat_{n_i}(\bar \kk)$,
  and thus, $w=\sum \la_i^{-1} n_i 1_i$, where $1_i$ is the identity matrix in
  $\Mat_{n_i}(\bar \kk)$. Therefore, $w$ is invertible in $A_{\bar \kk}$
  and thus in $A$.
\end{proof}
\begin{remark}
  An alternative proof which does not rely on the classification of
  semisimple algebras is given in \ocite{lauda-pfeiffer}*{Theorem 2.14}.
\end{remark}

Now we can return to \thref{t:adjunction-frob2}.

\begin{proof}[Proof of \thref{t:adjunction-frob2}]
  One direction is easy. Assume that we have adjunctions $\ev_A\dashv\coev_A$,
  $\coev_A\dashv\ev_A$. By \thref{t:adjunction-frob}, the first of these
  adjunctions defines on $A$ a structure of a symmetric Frobenius algebra. By
  \thref{t:adjunction-frob3}, the second adjunction gives an element $\sum
  \tilde x_i \otimes C\tilde x^i$ which is a separability idempotent (see
  \leref{l:separable}); thus, $A$ is separable.

  Conversely, let $A$ be a separable symmetric Frobenius algebra. By
  \thref{t:adjunction-frob}, this gives rise to adjunction $\ev_A\dashv\coev_A$,
  so we only need to construct the second adjunction $\coev_A\dashv\ev_A$. To do
  it, let $C=[w^{-1}]\in A/[A,A]$, where $w$ is the Euler element (see
  \thref{t:separable-frobenius}) and $\tilde e=e =\sum x_i \otimes x^i$; by
  \thref{t:adjunction-frob3} they define an adjunction $\coev_A\dashv
  \ev_A$.
\end{proof}
\begin{remark}
  The careful reader can note that the construction of bi-adjunction
  $\ev\dashv\coev$, $\coev\dashv\ev$ from a symmetric Frobenius algebra has an
  extra feature: the elements $e$, $\tilde e$, which are parts of these
  adjunctions, actually coincide. It is natural to ask whether it is always so:
  does having two adjunctions as above automatically imply that $\tilde e=e$?

  The answer is ``no'': in general, $\tilde e\ne e$; however, if both
  adjunctions exist, then we can modify data for one of them so that the equality
  $\tilde e=e$ is satisfied. See the discussion of ``cusp flip relation'' at the end
  of \seref{s:2d-xtqft-classification}.
\end{remark}


For future use, we also give here some further properties of separable symmetric
Frobenius algebras.

\begin{exercise}\label{ec:standard-frobenius}
  Let $A$ be a finite-dimensional semisimple algebra. Define the counit
  $\eps\colon A\to \kk$ by
  $$
  \eps(a) = \tr(L_a)
  $$
  where $L_a\colon A\to A$ is the operator of left multiplication by $a$.
  Prove that this defines on $A$ a structure of a symmetric Frobenius algebra,
  and for this structure, the Euler element \eqref{e:euler2} is $w=1$.

  (Hint: it suffices to consider the case $A=\Mat_n(\kk)$. Alternatively, to
prove that $(\, ,\, )$ is non-degenerate, you can prove that the kernel
of this bilinear form is a 2-sided ideal in $A$ which consists of
nilpotent elements, and thus for a semisimple algebra must be trivial.)
\end{exercise}



\begin{corollary}\label{c:euler3}
  Let $A$ be a separable symmetric Frobenius algebra, and let $w=\sum x_ix^i$
  be the Euler element defined by \eqref{e:euler2}. Then $w^{-1}$ is central,
  and the map
  \begin{equation}\label{e:center-isom2}
    A/[A,A] \to z(A)\colon [a]\mapsto \sum x_iax^i
  \end{equation}
  is an isomorphism; its inverse is given by
  \begin{equation}\label{e:center-isom3}
    z(A) \to A/[A,A]\colon z \mapsto w^{-1}z.
  \end{equation}
  In particular, $\sum x_i w^{-1}x^i=1$ and the
  element
  $$
    E= \sum w^{-1} x_i \otimes x^i = \sum x_i \otimes x^i w^{-1}
  $$
  is a symmetric separability idempotent.
\end{corollary}

Another reformulation of this corollary is the statement that the map
\begin{equation}\label{e:center-projector}
  \begin{aligned}
    p\colon A&\to A\\
            a&\mapsto w^{-1}\sum x_i a x^i
  \end{aligned}
\end{equation}
is a projector onto the center of $A$, and $\ker(p)=[A,A]$. Using the graphical
representation of operations in a Frobenius algebra as in \chref{c:frobenius},
this projector can be represented by \firef{f:frob-center}.

\begin{figure}[ht]
  \begin{tikzpicture}
    \node[dotnode] (om) at (0,0) {};
    \draw (om) arc(-90:270:0.4cm);
    \draw (om)--+(0,-0.5);
    \draw (om)--+(0, 1.5);
    \node[euler] at (0, -0.25){};
    %%%%%%%%%%%%
    \node at (0.75,0.5) {$=$};
    \node[dotnode] (om2) at (1.5,0) {};
    \draw (om2)--+(0,-0.5);
    \node[dotnode] (a1) at (1.25, 0.75){};
    \coordinate (a2) at (1.75, 0.75);
    \coordinate (b1) at (1.25, 0.25);
    \coordinate (b2) at (1.75, 0.25);
    \draw (om2) to[out=45,in=-90] (b2)
    to[out=90,in=-90] (a1)
    arc(180:0:0.25cm)
    to[out=-90,in=90] (b1)
    to[out=-90,in=135] (om2);
    \draw (a1)--+(0,0.75);
    \node[euler] at (1.5, -0.25){};
    %%%%%%%%%%%%
    \node at (2.25,0.5) {$=$};
    \node[dotnode] (om3) at (3,0) {};
    \draw (om3)--+(0,-0.5);
    \node[dotnode] (om4) at (3,1) {};
    \coordinate (a3) at (2.75, 0.75){};
    \coordinate (a4) at (3.25, 0.75);
    \coordinate (b3) at (2.75, 0.25);
    \coordinate (b4) at (3.25, 0.25);
    \draw (om3) to[out=45,in=-90] (b4)
    to[out=90,in=-90] (a3)
    to[out=90,in=-135] (om4);
    \draw (om4)
    to[out=-45,in=90] (a4)
    to[out=-90,in=90] (b3)
    to[out=-90,in=135] (om3);
    \draw (om4)--+(0,0.5);
    \node[euler] at (3, -0.25){};
  \end{tikzpicture}
  \caption{Projector onto $z(A)$, see \eqref{e:center-projector}. Red dot stands
  for operator of multiplication by $w^{-1}$.}
  \label{f:frob-center}
\end{figure}


%\begin{theorem}
%  Let $A, B$ be separable symmetric Frobenius algebras, and let ${}_BM_A$ be a
%  finite-dimensional $B-A$ bimodule. Let $M^*$ be the dual vector space,
%  considered as an $A-B$-bimodule (see \ecref{ec:dual-module-frob}).
%  Then $M^*$ is both left and right adjoint of $M$, considered as a
%  1-morphism in $\Alg$.
%\end{theorem}
%\begin{proof}
%  It suffices to prove that $M^*$ is right adjoint of $M$; then by the same
%  argument,  $M^{**} \simeq M$ is the right adjoint of $M^{*}$, so $M$
%  and $M^*$ are both left and right adjoints of each other.
%
%  Since $B$ is semisimple, any finite-dimensional left $B$-module $M$ is
%  projective and finitely generated. By \thref{t:alg-adjoint}, it has a right
%  adjoint, given by $M^\vee= \Hom_B(M, B)$. On the other hand, we have proved
%  before that for a symmetric Frobenius algebra, we have an isomorphism
%  $M^\vee\simeq M^*$ (see \ecref{ec:dual-module-frob}). Thus, $M^*$ is the
%  right adjoint of $M$.
%
%  It is instructive to give an explicit construction of the unit and counit of
%  adjunction, i.e. morphisms $A\to M^* \otimes_B M$, $M\otimes_A M^*\to B$. The
%  unit is given by $1\mapsto \sum   m^i\otimes m_i$, where  $m_i$ is a basis in
%  $M$ and $m^i$ the dual basis in $M^*$. To define the counit, let $b_i, b^i\in
%  B$ be such that $\sum b_i \otimes b^i$ is symmetric, bicentral, and  $\sum
%  b_ib^i=1$; by \coref{c:euler3}, such an element exists.  Define the counit
%map
%  $M\otimes_A M^*\to B$ by
%  $$
%    \eps(m\otimes m^*)=\sum_j \<b_jm, m^*\>b^j=\sum_j \<m, m^*b_j\>b^j
%  $$
%  where $\<\, , \,\>$ stands for the pairing $M\otimes M^*\to \kk$.
%
%  We leave it to the reader to check that so defined unit and counit satisfy
%  the adjunction equations, i.e. for any $m\in M$, $m^*\in M^*$ we have
%  \begin{equation}
%    \begin{aligned}
%      \sum \eps(m\otimes m^i)\cdot m_i&=m,\\
%      \sum m^i\cdot \eps(m_i\otimes  m^*)&=m^*.
%    \end{aligned}
%  \end{equation}
%\end{proof}
%
%
%%\begin{remark}
%%  One can notice that $\ssfrob$ is exactly the same as the category
%%  $\ssAlg$ of finite-dimensional semisimple algebras and bimodules: indeed,
%%  by \ecref{ec:standard-frobenius}, every finite-dimensional semisimple
%%  algebra has a symmetric Frobenius structure, and morphisms in $\ssfrob$
%%  do not use the Frobenius structure at all. So one might ask why bother
%%  introducing $\ssfrob$?
%%
%%  The correct answer is that you can think of $\ssfrob$ as an additional
%%  structure on the category $\ssAlg$: namely, choosing a Frobenius
%%  structure gives us a way of constructing an isomorphism between left and
%%  right adjoints of the same 1-morphism $M$. We will get back to it later.
%%\end{remark}



\section{Invariants and coinvariants}\label{s:invariants}
\coref{c:euler3}, which shows that for a separable symmetric Frobenius
algebra we have a natural isomorphism $A/[A,A]\simeq z(A)$, is a special
case of a more general result which we state here for future use.

Namely, let $M$ be a bimodule over an associative algebra $A$ (no further
restrictions yet). Define the spaces of invariants and coinvariants by
\begin{equation}\label{e:coinv}
  \begin{aligned}
    M^{A}&=\Hom(A,M)=\{m\in M\st am=ma \text{ for any } a\in A\}\\
    M_{A}&=A\otimes_{A^e} M = M/\<am-ma\>
  \end{aligned}
\end{equation}
where $\Hom$ is computed in the category of $A$-bimodules, or
equivalently, modules over $A^e=A\otimes A^\op$.

Standard results of algebra show that the functor of invariants is left
exact, and the functor of coinvariants is right exact. The corresponding
derived functors are called \emph{Hochschild (co)homology}\index{Hochschild
cohomology} of $A$ with
coefficients in $M$ (see, e.g., \ocite{pierce}*{Section 11.1}):
\begin{align*}
  HH^n(A,M)&=\Ext_{A^e}^n(A,M),\\
  HH_n(A,M)&=\Tor^{A^e}_n(A,M),
\end{align*}
so that $M^{A}=HH^0(A,M)$, $M_{A}=HH_0(A,M)$.
For example, for $M=A$, we get
\begin{align*}
  HH^0(A)&=A^{(A)}=z(A),\\
  HH_0(A)&=A_{(A)}= A/[A,A].
\end{align*}

For a general associative algebra, there is no relation between the functors of
invariants and coinvariants. However, if $A$ is a symmetric separable Frobenius
algebra, things are simplified. Namely, consider the map $p\colon M\to M$ given
by
$$
  m\mapsto \sum w^{-1} x_imx^i=\sum x_imx^i w^{-1}.
$$
As in \eqref{e:coinv-inv}, it is easy to see that it descends to a map
$M_{A}\to M^{A}$.
\begin{lemma}\label{l:coinv-separable}
  If $M$ is a bimodule over a separable symmetric Frobenius algebra $A$, then
  the map
  $$
  m \mapsto \sum x_imx^i
  $$
  is an isomorphism $M_{A}\isoto M^{A}$.
  \end{lemma}
For $M=A$, this is exactly the statement of \coref{c:euler3}. We leave it
to the reader to modify the proof of \coref{c:euler3} to this more general
case.


Note that since the functor of invariants is left exact and the functor of
coinvariants is right exact, this shows that for a separable symmetric Frobenius
algebra, both functors are exact so that higher Hochschild (co)homology vanishes.


As an application, given a right $A$-module $M_A$ and a left $A$-module
${}_A N$, there are two versions of tensor product over $A$:
\begin{equation}\label{e:tensor-types}
  \begin{aligned}
    M\otimes_A N &= \Coinv (M\otimes N) =
    M \otimes N/\<ma\otimes n- m\otimes an\>,\\
    M\otimes^A N &=\Inv(M\otimes N) \\
    &=\bigl\{x=\sum m_i\otimes n_i\in M\otimes N
    \st \sum m_ia\otimes n_i=\sum m_i\otimes an_i\bigr\}.
  \end{aligned}
\end{equation}
The first of these spaces (which is by far the more common) is a quotient of
$M\otimes N$; the second is a subspace. As an immediate corollary of
\leref{l:coinv-separable}, we see that for a separable symmetric Frobenius
algebra, we have a canonical isomorphism $M \otimes_A N \simeq M\otimes^A
N$.


